\section{Degradation Models}
\label{sec:degradations}

Degradations are the starting point of the restoration pipeline: they transform clean images into corrupted versions that models must learn to restore. Three types of parametric, reproducible degradations were implemented.

\subsection{Gaussian Noise}

Gaussian noise is the primary degradation used in our experiments. Given a pixel value $x$, the corrupted image is obtained as:
\begin{equation}
\tilde{x} = \text{clip}\!\left(x + \mathcal{N}(0, \sigma^2),\; 0,\; 255\right)
\label{eq:gaussian}
\end{equation}
where $\sigma$ controls noise intensity. We use $\sigma = 100$, a deliberately high value that severely corrupts image content and challenges the reconstruction capacity of all models. For reproducibility, each image receives a unique seed computed as $\texttt{seed} + \texttt{image\_index}$.

\subsection{Quantization and Dithering}

Quantization reduces the color depth of an image, mapping the original 256 levels per channel to $L = 2^b$ levels where $b$ is the number of bits per channel.

\subsubsection{Simple Quantization}
Direct depth reduction:
\begin{equation}
q(x) = \text{round}\!\left(\frac{x}{255} \cdot (L-1)\right) \cdot \frac{255}{L-1}
\end{equation}

\subsubsection{Random Dithering}
Before quantization, uniform noise is added in the interval $[-s/2, s/2]$ where $s = 255/(L-1)$ is the quantization step:
\begin{equation}
\tilde{x} = q\!\left(x + \mathcal{U}\!\left(-\frac{s}{2},\; \frac{s}{2}\right) \cdot \text{noise\_strength}\right)
\end{equation}
This breaks the banding patterns typical of simple quantization.

\subsubsection{Ordered Dithering (Bayer)}
Uses Bayer matrices ($2\!\times\!2$, $4\!\times\!4$, or $8\!\times\!8$) as spatially-varying thresholds. The Bayer matrix is normalized to $[-0.5, 0.5]$ and added to the image before quantization, producing regular, visually pleasing patterns.

\subsubsection{Floyd-Steinberg Dithering}
An error diffusion algorithm~\cite{floyd} that processes the image pixel by pixel. The quantization error $e = x - q(x)$ is distributed to adjacent unprocessed pixels according to the kernel:
\begin{equation}
\begin{bmatrix}
\cdot & \cdot & \cdot \\
\cdot & \ast & \frac{7}{16} \\[3pt]
\frac{3}{16} & \frac{5}{16} & \frac{1}{16}
\end{bmatrix}
\end{equation}
where $\ast$ denotes the current pixel. This produces perceptually superior results at the cost of sequential processing.

\subsection{Salt-and-Pepper Noise}

Salt-and-pepper noise replaces a fraction $p$ of pixels with extreme values: the parameter \texttt{salt\_vs\_pepper} controls the ratio between white pixels (255, \textit{salt}) and black pixels (0, \textit{pepper}).

\subsection{Automated Dataset Generation}

The module \texttt{generate\_degraded\_dataset.py} provides batch processing of entire image directories with automatic path generation following a consistent naming convention (\eg, \texttt{data/\allowbreak{}degraded/\allowbreak{}gaussian/\allowbreak{}sigma\_100/\allowbreak{}DIV2K\_train\_HR/}). Generation is skipped if the degraded dataset already exists, ensuring reproducibility without redundant computation.
